% ---
% id: meat-core/req:exact-admission
% kind: req
% title: "Exact admission of declared data"
% summary: >-
%   The engine admits a declaration only if it is exactly admissible and
%   refuses every other, with no tolerance on declared data, save a reference
%   companion, admitted on its own terms.
%   [[prop:optol-nobound]] shows, for the measurement on Q, that closeness to
%   admissible data could license no key-length claim.
% prominence: standard
% status: implemented
% uses: [def:declaration, prop:optol-nobound]
% mentions: [foundations/req:declared-tolerance, foundations/req:consumer-decides-exactly, req:companion-admission, meat-core/frank-wolfe/rem:fo-arithmetic]
% implements: certiqs_key.meat_core.declaration.general_validation:validate_general_declaration
% assumes: []
% refs: []
% provenance: newly-derived
% lean:
% ---

Exact admissibility asks the conditions of
Definition~\ref{def:declaration} of the declared
\tref{binary64 numbers} themselves, not of reals they approximate. For a
declaration with one Alice sector, one \tref{measured factor} and Alice's
measurement in \tref{POVM form}, this means that, block by block, each
declared operator is exactly Hermitian, each measurement operator and the
declared $\tau_A$ are exactly positive semidefinite, and each completeness
sum, formed in the \tref{declared order} \eqref{eq:declorder}, is exactly
the identity with no addition rounding; the paragraph of that definition on
exact admissibility in the general format says how the conditions read
otherwise. How the rule meets an instance that moves its own data within
its \tref{declared tolerance} $\tnum$ is
\xref{foundations/req:declared-tolerance}, read with
\xref{foundations/req:consumer-decides-exactly}.

\begin{requirement}\label{req:exact-admission}
The engine takes no tolerance on \tref{declared data} and decides no
condition on it more loosely than exactly
(\xref{foundations/req:consumer-decides-exactly}): it admits a declaration
only if the declaration is \tref{exactly admissible} in the sense of
Definition~\ref{def:declaration}, and refuses it otherwise. No
perturbation of a declared family is licensed to carry the key length of
the family it approximates. Proposition~\ref{prop:optol-nobound} is the
reason: it shows, for the measurement on $\Qsys$, that no tolerance could
license a key-length claim. A compiled reference companion is not admitted
under this requirement but under Requirement~\ref{req:companion-admission},
which decides exactly what it shares with the declaration and decides none
of its family; the binary64 numbers it is evaluated at are no declared
family whose key length is claimed, since the bound evaluated there is that
of the reference, at the unbudgeted arithmetic of
Remark~\ref{rem:fo-arithmetic}.
\end{requirement}
